\begin{theorem}[NestedShrinkInterval Cantor downstream route]
\label{thm:nested-shrink-interval-cantor-downstream-route}
Every $\CantorDiagonalCompactnessUp$ or $\CantorIntersectionUp$ consumer of a $\NestedShrinkIntervalUp$ packet reads the nested source only through nested closed intervals, located endpoint rows, $\DyadicRatCoreUp$ diameter witnesses, $\StreamNameUp$ finite windows, $\RegSeqRatUp$ regular readback, and the terminal $\RealUp$ seal. The downstream route exports a Cantor diagonal or intersection-facing interval certificate, but it does not select a limit point, assert ambient completeness, identify quotient-real names, or add a host compactness principle outside the displayed rows.
\end{theorem}
\begin{proof}
Project the nested-shrink carrier through the child routes named by the hub. The compact-interval and rational-locator rows provide nested closed intervals and located endpoints, while the dyadic and StreamName rows expose only finite diameter and window observations.

The Cantor-facing consumer then passes through the RegSeqRat and Real rows already used by the completion and sibling routes. Since these rows carry finite readback and terminal naming rather than a chosen point, the consumer obtains only the displayed interval certificate.

No carrier coordinate names a selected limit point, host completeness theorem, quotient equality, or ambient compactness object. Thus $\CantorDiagonalCompactnessUp$ and $\CantorIntersectionUp$ consume $\NestedShrinkIntervalUp$ only through the finite interval route stated above.
\end{proof}
